\documentclass[UTF8, 11pt]{ctexart}

\usepackage[margin=1in]{geometry}
\usepackage{amsmath, amssymb}
\usepackage{booktabs}
\usepackage{hyperref}
\usepackage{url}
\usepackage{enumitem}
\setlist{leftmargin=*}
\setlength{\parskip}{0.35em}
\setlength{\parindent}{0pt}

\title{\textbf{阶段 1：背景与相关工作}\\
面向通信高效分布式优化的可扩展图划分算法}
\author{小组成员：姓名 1，姓名 2，姓名 3}
\date{CS240：算法设计与分析，2026 年春季学期}

\begin{document}

\maketitle

\section{课内算法基础}

在介绍具体应用之前，有必要先回顾算法设计与分析中的几个基础概念。许多现代
数据问题都需要先选择合适的数学表示，才能变成可分析、可实现的算法问题。本项
目使用的表示是图。图把一个系统抽象为顶点和边：顶点表示对象或计算单元，边
表示关系、依赖或通信链路。如果边带有权重，权重通常表示代价、相似度、容量或
交互强度。

图问题是算法课中的核心内容，因为它把建模、优化和复杂度联系在一起。最短路、
最小生成树、匹配、流、割和图划分，都是把实际应用转化为图模型之后得到的算法
问题。其中，割问题与本项目最相关。一个割把图中的顶点分成不同集合，并衡量
跨越这些集合的边。普通最小割在某些设置下可以被精确高效求解，但如果再加入
“每个部分大小要均衡”这样的额外约束，问题就会明显变难。

因此，本项目从一个熟悉的课内思想出发：先把问题表示成图，再研究一个更困难的
割问题，即平衡图划分。这里涉及的主要算法主题包括：
\begin{itemize}
    \item \textbf{目标函数：} 最小化跨分区边权重；
    \item \textbf{约束条件：} 保持分区均衡，避免某个 worker 承担过多任务；
    \item \textbf{启发式算法：} 当精确优化代价过高时，使用实际有效的方法；
    \item \textbf{复杂度与可扩展性：} 分析图规模增大时运行时间如何变化。
\end{itemize}

有了这些背景，五个阶段就可以自然地读成一个递进过程。阶段 1 建立图模型并介绍
相关工作；阶段 2 选择并分析代表性算法；后续阶段再实现算法、测试算法，并讨论
可能的拓展。

\section{应用背景}

本项目研究经典图划分算法在通信高效分布式优化中的应用。在现代机器学习、科学
计算和大规模数据处理中，许多问题规模过大，无法只依靠单台机器完成。常见做法
是将任务分配到多个 worker 或智能体上，每个 worker 负责一部分数据、变量或计算
任务，并通过消息传递来协调局部更新。

在这类系统中，通信经常是主要瓶颈之一。局部数值计算通常比较便宜，而同步、
跨机器数据传输和重复消息交换可能非常昂贵。如果两个依赖关系很强的任务被分配
到不同 worker 上，它们之间的交互就会变成跨 worker 通信。如果大量依赖边跨越
worker 边界，分布式算法就可能把大量时间花在传输信息上，而不是花在真正的计算
上。

图划分为这一问题提供了自然的算法抽象。我们将计算过程建模为图
$G=(V,E,w)$。顶点表示数据块、变量块、计算任务或智能体；边表示依赖关系、消息
传递需求或通信链路；边权表示通信强度或通信代价。将顶点集合 $V$ 划分为 $k$ 个
平衡部分，就对应于把任务分配给 $k$ 个 worker。理想划分应当同时满足两个目标：
保持每个 worker 的负载接近，并尽量减少跨分区边的总权重。

\section{算法形式化}

本项目的核心计算问题是平衡 $k$ 路图划分。给定一个无向加权图 $G=(V,E,w)$ 和
整数 $k$，目标是找到互不相交的集合 $V_1,\ldots,V_k$，使得
\[
    V = V_1 \cup \cdots \cup V_k,
    \qquad V_i \cap V_j = \emptyset \text{ for } i\neq j.
\]
割边目标函数定义为
\[
    \mathrm{cut}(V_1,\ldots,V_k)
    =
    \sum_{(u,v)\in E,\; u\in V_i,\; v\in V_j,\; i\neq j} w_{uv}.
\]
为了避免某一个 worker 被分配到过多顶点，还需要加入负载均衡约束，例如
\[
    |V_i| \leq (1+\epsilon)\frac{|V|}{k},
    \qquad i=1,\ldots,k.
\]
优化目标是在满足均衡约束的前提下最小化割边权重。

这一形式化与许多经典算法直接相关。最小割和最大流可以解决某些精确割问题，但
平衡图划分更困难，因为均衡约束排除了把全部顶点放入同一侧的平凡解。谱方法用
图拉普拉斯矩阵的特征向量寻找低割划分。Kernighan--Lin 等局部搜索方法通过交换
顶点迭代改进划分。METIS 等多层方法则通过图粗化、粗图划分和细化来提升大规模
图上的可扩展性。

\section{代表性文献}

\textbf{Kernighan--Lin 算法。}
Kernighan 和 Lin 提出了经典的图划分启发式算法。该方法从一个初始二分开始，
反复选择能够降低割值的顶点交换。虽然原论文主要研究图二分，而不是一般的
$k$ 路划分，但其局部搜索和增益最大化思想仍然非常重要。本项目复现这一细化
思想，并通过在不同分区对之间执行成对 Kernighan--Lin 二分，将其改写为实践中的
$k$ 路划分细化过程。

\textbf{谱图划分。}
谱划分使用图拉普拉斯矩阵 $L=D-A$，其中 $D$ 是度矩阵，$A$ 是邻接矩阵。对应于
第二小特征值的特征向量称为 Fiedler 向量，它可以看作某种图割问题的连续松弛
解。按照 Fiedler 向量排序并二分顶点，可以得到一个自然的图二分；递归二分可以
进一步得到多个分区。谱方法的重要性在于，它揭示了图割、线性代数和连续优化
松弛之间的联系。

\textbf{多层图划分。}
Karypis 和 Kumar 的 METIS 风格多层框架显著提升了实际图划分算法的速度和质量。
该框架首先把图逐步粗化成更小的图，在最粗图上求划分，然后再逐层投影回原图并
进行细化。虽然本项目没有完整实现多层框架，但 METIS 是重要参照，因为它说明了
局部细化如何与粗化结合，从而适用于大规模不规则图。

\textbf{网络上的分布式优化。}
分布式次梯度法、一致性算法和去中心化梯度方法研究多个智能体如何在通信网络上
协同优化。在这些方法中，通信图决定哪些智能体可以交换信息。本项目把这一领域
作为现代应用背景：图划分不仅是图处理任务，也可以作为减少分布式算法跨 worker
通信的工具。

\section{选题的算法研究价值}

该选题具有明显的算法研究价值。第一，平衡图划分同时包含目标函数和硬约束。
割值很小但极不均衡的划分无法用于分布式计算；完全均衡但跨边很多的划分也会
带来巨大通信代价。算法必须在割边代价和负载均衡之间做权衡。

第二，该问题能展示精确算法、连续松弛和启发式算法之间的差异。本项目比较贪心
均衡基线、谱划分和 Kernighan--Lin 局部细化，分别代表了不同的算法设计思想。

第三，应用场景具有现代性。图划分出现在分布式机器学习、图神经网络训练、稀疏
矩阵计算和多智能体优化中。经典算法工具因此可以直接服务于当前的数据和机器
学习系统问题。

第四，该问题可以通过实验进行量化评估。我们可以测量割边权重、归一化割、负载
均衡比例、运行时间、可扩展性，以及一致性仿真中的跨分区通信量。这些指标把
形式化目标与可观察的系统行为联系起来。

\section{项目定位}

本项目选择课程参考选题中的 \textbf{Topic 2: Community Detection / Graph
Partitioning}。与常见的社交网络社区发现不同，我们把重点放在分布式优化中的
图划分。该选择仍然完全符合参考选题，因为核心算法问题是图划分；同时，它也把
经典问题连接到了现代机器学习和数据系统应用。

本项目主要复现 Kernighan--Lin 图划分启发式算法，并实现谱二分和贪心均衡划分
作为基线。预期结果是展示经典图算法可以在保持负载均衡的同时降低跨 worker
通信成本。

\begin{thebibliography}{9}

\bibitem{kernighan1970}
B. W. Kernighan and S. Lin.
\newblock An efficient heuristic procedure for partitioning graphs.
\newblock \emph{Bell System Technical Journal}, 1970.

\bibitem{karypis1998}
G. Karypis and V. Kumar.
\newblock A fast and high quality multilevel scheme for partitioning irregular graphs.
\newblock \emph{SIAM Journal on Scientific Computing}, 1998.

\bibitem{nedic2009}
A. Nedic and A. Ozdaglar.
\newblock Distributed subgradient methods for multi-agent optimization.
\newblock \emph{IEEE Transactions on Automatic Control}, 2009.

\end{thebibliography}

\end{document}
